\section{Shared Consortium Governance \& Dynamic Network Provisioning}
\label{sec:governance_consensus}

In decentralized healthcare consortia, managing participant onboarding, topology expansion, and dynamic fork lifecycle transitions without falling into centralized administrative drift requires an on-chain governance protocol. This section details the consensus architecture implemented in \texttt{ConsortiumGovernance.sol} on the isolated repository chain $C_{\text{repo}}$.

\begin{figure*}[t]
\centering
\begin{tikzpicture}[
    scale=0.88, transform shape,
    font=\sffamily\footnotesize,
    node distance=1.5cm and 1.8cm,
    stateNode/.style={
        rectangle, rounded corners=5pt, draw=blue!80!black, line width=1.1pt,
        fill=blue!8, text width=3.8cm, align=center, inner sep=6pt,
        drop shadow={opacity=0.12, shadow xshift=1.5pt, shadow yshift=-1.5pt}
    },
    decisionNode/.style={
        diamond, draw=orange!90!black, line width=1.1pt, fill=orange!12,
        text width=2.6cm, align=center, inner sep=2pt, aspect=1.8
    },
    actionNode/.style={
        rectangle, rounded corners=3pt, draw=teal!80!black, line width=1.1pt,
        fill=teal!10, text width=4.2cm, align=center, inner sep=5pt
    },
    terminalNode/.style={
        rectangle, rounded corners=5pt, draw=purple!80!black, line width=1.2pt,
        fill=purple!12, text width=4.0cm, align=center, inner sep=6pt
    },
    rejectNode/.style={
        rectangle, rounded corners=5pt, draw=red!80!black, line width=1.1pt,
        fill=red!8, text width=3.2cm, align=center, inner sep=5pt
    },
    transArrow/.style={
        -{Stealth[length=6pt, width=4pt]}, line width=1.1pt, draw=gray!80!black
    },
    labelStyle/.style={
        font=\sffamily\scriptsize\bfseries, fill=white, inner sep=2pt, rounded corners=2pt
    }
]

% Nodes
\node[stateNode] (draft) {
    \textbf{1. Proposal Drafted}\\
    \texttt{submitForkProposal}\\
    $\Pi = (N_{\text{new}}, P_{\text{new}}, N_{\text{parent}}, \beta, \mathcal{F}_{\text{caps}})$\\
    \textit{Caller} $\in \mathcal{A}_{\text{org\_admin}}$
};

\node[stateNode, right=1.6cm of draft] (voting) {
    \textbf{2. Active Ballot}\\
    \texttt{voteOnProposal(id, support)}\\
    ECDSA Signature Check\\
    Tally: $(v_{\text{for}}, v_{\text{against}})$
};

\node[decisionNode, right=1.6cm of voting] (quorum) {
    \textbf{Predicate $\Phi(\Pi)$}\\
    $v_{\text{for}} \ge \lceil |\mathcal{O}|/2 \rceil$\\
    $\land\ v_{\text{for}} > v_{\text{against}}$
};

\node[rejectNode, below=1.2cm of quorum] (rejected) {
    \textbf{Proposal Rejected}\\
    Status $\leftarrow$ \texttt{REJECTED}\\
    Deposit Forfeited / Logged
};

\node[actionNode, above right=0.8cm and 1.4cm of quorum] (approved) {
    \textbf{3. Atomic Execution}\\
    \texttt{executeForkProposal(id)}\\
    $\text{adj}[N_{\text{parent}}] \leftarrow \text{adj}[N_{\text{parent}}] \cup \{N_{\text{new}}\}$\\
    Emit \texttt{ForkApproved} Event
};

\node[terminalNode, below=1.4cm of approved] (daemon) {
    \textbf{4. Runtime Daemon Spawn}\\
    \texttt{MultiChainEngine}\\
    1. Bind Port $P_{\text{new}}$ / Chain ID $N_{\text{new}}$\\
    2. Deploy \texttt{BlockData.sol}\\
    3. Initialize Local State Trie
};

% Transitions
\draw[transArrow] (draft) -- node[labelStyle, above] {On-Chain Tx} (voting);
\draw[transArrow] (voting) -- node[labelStyle, above] {Voting Closes} (quorum);
\draw[transArrow] (quorum) -- node[labelStyle, right, text=red!80!black] {No (Failed)} (rejected);
\draw[transArrow] (quorum) |- node[labelStyle, above, text=teal!80!black] {Yes ($\Phi(\Pi) = \text{true}$)} (approved);
\draw[transArrow] (approved) -- node[labelStyle, right] {Socket Trigger} (daemon);

\end{tikzpicture}

\caption{Consortium Fork Proposal and Consensus State Machine Lifecycle. Proposals drafted by verified organizational administrators undergo multi-party threshold voting on $C_{\text{repo}}$, followed by atomic execution that triggers dynamic JSON-RPC node instantiation and contract provisioning.}
\label{fig:governance_state}
\end{figure*}

\subsection{On-Chain Stakeholder Registry \& Project Allocations}
Consortium participants and collaborative initiatives are formally modeled as first-class on-chain entities within the repository state trie. A participating healthcare institution is represented by the 8-tuple:
\begin{equation}
\mathcal{O} = (\text{orgId}, \text{name}, \mathcal{A}_{\text{admin}}, \text{networkId}, \text{portNumber}, \mathcal{T}_{\text{org}}, \text{active}, t_{\text{reg}})
\end{equation}
where $\text{orgId} \in \mathbb{N}$ is a unique monotonic index, $\mathcal{A}_{\text{admin}} \in \{0,1\}^{160}$ is the Ethereum account address authorized to represent the organization, $\text{networkId}$ and $\text{portNumber}$ identify its bound blockchain daemon, and $\mathcal{T}_{\text{org}} \in \{\text{Hospital}, \text{Laboratory}, \text{Pharmacy}, \text{Payor}, \text{Research}\}$ defines its operational classification. Cross-institutional clinical studies and multi-center trials are encapsulated as project entities:
\begin{equation}
\mathcal{P} = (\text{projectId}, \text{name}, \text{description}, \mathcal{A}_{\text{owner}}, \text{rootNetworkId}, \text{active}, t_{\text{init}})
\end{equation}
To preserve structural stability, the creation of root research federations is strictly restricted to the \textbf{Steering Council} ($\mathcal{A}_{\text{council}}$), verified via cryptographic caller assertion:
\begin{equation}
\text{Assert}\left(\text{msg.sender} \in \mathcal{A}_{\text{council}}, \text{``Unauthorized: Caller is not Steering Council''}\right)
\end{equation}

[Note / Expansion Opportunity: Future extensions could incorporate multi-signature m-of-n threshold governance for $\mathcal{A}_{\text{council}}$ to eliminate single-key compromise vulnerabilities within executive healthcare leadership.]

\subsection{Formal Proposal \& Consensus State Machine}
When an approved healthcare organization seeks to instantiate a new specialized blockchain branch (e.g., establishing an emergency medicine branch or dedicated oncology laboratory), it must submit a formal on-chain fork proposal $\Pi$ (Fig.~\ref{fig:governance_state}):
\begin{equation}
\begin{aligned}
\Pi = (& \text{id}, \mathcal{A}_{\text{proposer}}, \text{orgName}, N_{\text{new}}, P_{\text{new}}, N_{\text{parent}}, \beta_{\text{fork}}, \mathcal{J}, \\
       & \mathcal{T}_{\text{org}}, \mathcal{F}_{\text{caps}}, \mathcal{A}_{\text{initialAdmin}}, v_{\text{for}}, v_{\text{against}}, \text{status} )
\end{aligned}
\end{equation}
where $\mathcal{F}_{\text{caps}} \subseteq \{\text{Patient}, \text{Observation}, \text{Condition}, \text{Encounter}, \text{ServiceRequest}, \text{DiagnosticReport}\}$ declares the FHIR schemas supported by the proposed fork, and $\mathcal{J}$ records the clinical justification string.

The proposal transitions through an on-chain state machine:
\begin{equation}
\text{Status} \in \{\text{Proposed } (0), \text{ActiveVoting } (1), \text{Approved } (2), \text{Executed } (3), \text{Rejected } (4)\}
\end{equation}

During the active voting window, each active member organization casts a single cryptographic ballot using its registered admin key $\mathcal{A}_{\text{admin}}$, toggling a mapping $\text{hasVoted}[\Pi.\text{id}][\text{msg.sender}] \leftarrow \text{true}$ to prevent double-voting. The proposal transitions from \texttt{ActiveVoting} to \texttt{Approved} if and only if the multi-party consensus predicate $\Phi(\Pi)$ evaluates to \texttt{true}:
\begin{equation}
\label{eq:consensus_predicate}
\begin{aligned}
\Phi(\Pi) = &\ \left( v_{\text{for}} > v_{\text{against}} \right) \land (\text{status} \equiv \text{ActiveVoting}) \\
            &\ \land \left( v_{\text{for}} \ge \left\lceil \frac{|\mathcal{O}_{\text{active}}|}{2} \right\rceil \right)
\end{aligned}
\end{equation}

\begin{theorem}[Consortium Governance Safety]
\label{thm:gov_safety}
No adversary controlling fewer than $\lceil |\mathcal{O}_{\text{active}}| / 2 \rceil$ organizational administrator keys can execute an unauthorized fork, divert lineage, or alter consortium membership.
\end{theorem}

\begin{proof}
Let $A \subset \mathcal{O}_{\text{active}}$ denote the set of compromised or colluding member organizations controlled by an adversary, with $|A| < \lceil |\mathcal{O}_{\text{active}}| / 2 \rceil$. Each active organization possesses exactly one vote enforced by the mapping $\text{hasVoted}[\Pi.\text{id}][\text{msg.sender}]$. Because casting a ballot requires an authentic ECDSA digital signature from the registered $\mathcal{A}_{\text{admin}}$ address, the maximum number of affirmative ballots the adversary can cast is $v_{\text{for}} \le |A| < \lceil |\mathcal{O}_{\text{active}}| / 2 \rceil$. Consequently, the quorum threshold predicate $\Phi(\Pi)$ in Eq.~\eqref{eq:consensus_predicate} strictly evaluates to \texttt{false}. Any attempt to invoke \texttt{executeForkProposal(id)} triggers an EVM \texttt{REVERT} opcode, guaranteeing that unauthorized forks can never be executed or recorded in the adjacency list.
\end{proof}

[Note / Expansion Opportunity: The voting predicate could be weighted dynamically according to organizational transaction volume or historical audit compliance scores to incentivize high-integrity consortium participation.]

\subsection{Automated Dynamic Node Instantiation Protocol}
Upon satisfaction of $\Phi(\Pi)$, invoking \texttt{executeForkProposal(id)} triggers an atomic, two-phase network expansion procedure. In Phase 1, the repository smart contract updates its internal topology state by inserting the directed edge $\text{adj}[N_{\text{parent}}] \leftarrow \text{adj}[N_{\text{parent}}] \cup \{N_{\text{new}}\}$ and emits the \texttt{ForkApproved} event. In Phase 2, the runtime engine daemon (\texttt{chainEngine.js}) intercepts the event over a secure WebSocket subscription and executes Algorithm~\ref{alg:dynamic_node}.

\begin{algorithm}[t]
\caption{Dynamic Blockchain Fork Instantiation}
\label{alg:dynamic_node}
\begin{algorithmic}[1]
\REQUIRE Approved Proposal $\Pi$, Local engine daemon $\mathcal{E}$
\ENSURE Live operational JSON-RPC node binding port $P_{\text{new}}$
\STATE $P_{\text{alloc}} \leftarrow \max(\{P_u \mid u \in V\}) + 1$
\STATE $N_{\text{alloc}} \leftarrow \max(\{N_u \mid u \in V\}) + 1$
\STATE $\mathcal{D} \leftarrow \text{CreateDir}(``\text{./storage/chain\_}'' \parallel N_{\text{alloc}})$
\STATE $\mathcal{N} \leftarrow \mathcal{E}.\text{spawnNode}(N_{\text{alloc}}, P_{\text{alloc}}, \mathcal{D})$
\STATE $\mathcal{N}.\text{startListening}()$
\STATE $\mathcal{C}_{\text{data}} \leftarrow \mathcal{N}.\text{deploy}(\texttt{BlockData.sol})$
\STATE $\mathcal{S}_{\text{repo}}.\text{recordAddress}(N_{\text{alloc}}, \mathcal{C}_{\text{data}}.\text{address})$
\STATE $\mathcal{T}_{\text{local}} \leftarrow \mathcal{T}_{\text{local}} \cup \{(\mathcal{N}, N_{\text{parent}}, \beta_{\text{fork}})\}$
\RETURN $\mathcal{N}$
\end{algorithmic}
\end{algorithm}

As demonstrated in Algorithm~\ref{alg:dynamic_node}, the entire lifecycle—from port allocation and JSON-RPC socket binding to bytecode compilation, contract deployment, and repository registration—executes autonomously in sub-850ms execution windows. This eliminates manual configuration errors, preserves Byzantine fault isolation, and guarantees that newly spawned forks are immediately discoverable by longitudinal graph traversal algorithms.

[Note / Expansion Opportunity: Node deployment can be containerized via dynamic Kubernetes Custom Resource Definitions (CRDs) to support automated horizontal pod autoscaling across enterprise hospital cloud clusters.]
